\section{结论}

本论文基于已有的 低压浮动反相放大器 (Low-Voltage Floating Inverter Amplifier, LVFIA)，鉴于其在低电压环境下增益不足、线性下降、共模漂移敏感等问题，
提出了一种基于 相关电平移位 (Correlated Level Shifting, CLS) 技术的创新架构设计方法。
通过理论分析、仿真验证和实验测试，系统评估了CLS在提升动态放大器性能方面的有效性，
并与传统固定偏置式电平移位方法进行了对比。结果表明，CLS能够显著提升动态放大器的增益和线性度，
同时降低共模漂移敏感性，为未来低电压模拟电路设计提供了新的思路和方法。

